\chapter{Organisme d'accueil et étude de l'existant}
\section{Contexte du stage}
AMALYTICS est l'organisme d'accueil du stage, situé à Paris. Dans le cadre de ce rapport, sa présentation est volontairement limitée à un contexte générique de travail autour de la donnée et de solutions numériques appliquées à la santé. Cette précaution permet de décrire le projet sans exposer d'information interne ou de donnée sensible.

Le stage s'inscrit dans un besoin opérationnel fréquent : transformer des informations cliniques non structurées en données structurées exploitables par un système de suivi d'étude. Le périmètre couvert par la solution concerne des fichiers locaux au format PDF, TXT ou Markdown. Les sorties visées sont des mises à jour traçables de champs eCRF et des exports compatibles avec des workflows d'évaluation ou des classeurs Excel.

\section{Problématique}
Un document médical n'est pas une base de données. Il mélange texte libre, tableaux, résultats numériques, unités, antécédents, dates et conclusions. Une même notion peut être exprimée avec plusieurs libellés. Dans un bilan biologique, un résultat peut apparaître avant son analyte, être séparé par des points de conduite ou partagé entre deux colonnes. Dans un compte rendu d'imagerie, une mesure de référence peut coexister avec une mesure courante, une lésion non cible et une conclusion RECIST.

Une approche monolithique consistant à envoyer tout le document à un unique modèle d'extraction ne fournit pas un contrôle suffisant sur le contexte, la provenance et la validation. Elle rend aussi l'analyse des erreurs difficile. La problématique retenue est donc la suivante : \emph{comment extraire de manière fiable, configurable et traçable les informations utiles d'un document médical hétérogène, tout en limitant les écritures non justifiées dans un eCRF ?}

\section{Étude de l'existant}
Les solutions existantes peuvent être regroupées selon trois familles. La première repose sur la saisie manuelle assistée par une interface eCRF. Elle offre un contrôle humain direct, mais consomme du temps et introduit un risque de transcription. La deuxième utilise de l'OCR ou des règles fixes. Elle est efficace sur des formats stables, mais fragile face aux variations de mise en page et de vocabulaire. La troisième combine recherche d'information et modèles de langage ou d'extraction. Elle peut traiter le texte libre, mais exige des garde-fous afin de conserver l'évidence source et de contrôler les valeurs générées.

\begin{table}[H]
\centering
\caption{Comparaison qualitative des principales approches}
\begin{tabularx}{\textwidth}{lXXX}
\toprule
Approche & Atout principal & Limite principale & Usage adapté\\
\midrule
Saisie manuelle & Contrôle explicite & Lenteur et répétition & Cas complexes ou exceptions\\
OCR + règles & Déterminisme & Sensibilité au format & Documents très homogènes\\
Extraction monolithique & Mise en œuvre rapide & Faible explicabilité & Prototypes limités\\
Pipeline modulaire & Traçabilité et testabilité & Conception plus exigeante & Production contrôlée\\
\bottomrule
\end{tabularx}
\end{table}

La solution conçue se place dans cette dernière catégorie. Elle utilise des règles déterministes lorsqu'une structure fiable est disponible, par exemple pour des lignes de laboratoire, et réserve les mécanismes d'extraction plus souples aux contenus narratifs. Les composants sont cohérents avec les possibilités des outils employés : \texttt{pypdf} pour une lecture légère des PDF \cite{pypdf}, Docling pour une analyse structurée optionnelle \cite{docling}, et Qdrant pour la recherche vectorielle hybride \cite{qdrant}.

\section{Besoins fonctionnels}
Les besoins retenus ont été formulés comme des capacités vérifiables :
\begin{itemize}
  \item charger localement un document et conserver ses métadonnées ;
  \item identifier le type de document et la famille de champs pertinente ;
  \item reconstruire les structures utiles, notamment les lignes biologiques ;
  \item extraire des valeurs, unités, temporalités et preuves textuelles ;
  \item normaliser les valeurs selon le schéma d'étude ;
  \item valider la compatibilité avec le champ cible avant toute exportation ;
  \item exporter les propositions sous forme JSON, CSV ou XLS ;
  \item produire des artefacts d'évaluation reproductibles.
\end{itemize}

\section{Exigences non fonctionnelles}
La sensibilité du domaine impose de privilégier la traçabilité, la sécurité des écritures et la reproductibilité. Chaque valeur doit conserver un lien avec le texte qui la justifie. Le système doit pouvoir isoler les index par étude, éviter de remplir des colonnes déjà renseignées selon une politique explicite et signaler les cas non validés. Les performances mesurées dans ce stage sont des métriques d'extraction hors ligne ; elles ne doivent pas être assimilées à des performances cliniques.

\subsection{Traçabilité et explicabilité}
La valeur d'une extraction ne peut pas être appréciée uniquement par un score numérique. Une proposition doit être accompagnée de l'extrait de texte dont elle est issue, de la position dans le document, de la stratégie qui l'a produite, de sa clé canonique et du résultat des validations appliquées. Cette information est particulièrement utile lorsqu'une valeur est rejetée : elle permet de déterminer si la cause se situe dans le parsing, la recherche, l'extraction, la normalisation ou la règle de sortie.

La solution conserve donc des objets intermédiaires au lieu de ne conserver que l'export final. Cette décision augmente le volume d'informations techniques, mais facilite les audits et le débogage. Elle rend également possible une revue humaine ciblée sur les valeurs incertaines, plutôt que sur l'ensemble du document.

\subsection{Configurabilité}
Les études cliniques ne partagent pas toujours les mêmes colonnes, unités, seuils ni temporalités. La configuration doit par conséquent être indépendante du code d'orchestration. Le schéma d'étude versionné est l'interface principale de cette configurabilité. Il permet de déclarer une nouvelle famille de champs, d'associer des synonymes et d'indiquer la stratégie d'extraction attendue sans modifier le cœur de la pipeline.

\subsection{Qualité et maintenabilité}
Le projet doit être exécutable par un autre développeur ou par une équipe de reprise. Les scripts de démonstration, les tests déterministes, les fichiers Docker Compose et la documentation de passation répondent à cette exigence. La modularité rend les responsabilités lisibles : une modification du mapping d'alias biologique ne nécessite pas de modifier la couche de recherche vectorielle. Cette séparation limite aussi l'effet de bord des corrections successives.

\section{Méthodologie de travail : CRISP-DM}
Afin de structurer le développement de la solution, nous avons adopté une démarche inspirée de CRISP-DM (\emph{Cross-Industry Standard Process for Data Mining}) \cite{crispdm}. Cette méthodologie organise un projet orienté données en six phases itératives : compréhension métier, compréhension des données, préparation des données, modélisation, évaluation et déploiement. Son caractère cyclique est particulièrement adapté au projet : les résultats d'évaluation ont permis d'identifier des erreurs de parsing ou de sélection de contexte, puis d'améliorer les composants concernés.

Dans notre cas, CRISP-DM a été adapté au développement d'une pipeline d'extraction d'informations médicales destinée à assister l'alimentation des eCRF. Il ne s'agit pas d'un projet de prédiction clinique : la phase de modélisation recouvre ici les stratégies de recherche, d'extraction, de normalisation et de validation des informations documentaires.

\begin{figure}[H]
\centering
\begin{tikzpicture}[every node/.style={draw,circle,minimum size=2.1cm,align=center,font=\scriptsize,fill=softblue}, >=Latex]
\node (business) at (90:3.0cm) {Compréhension\\métier};
\node (data) at (30:3.0cm) {Compréhension\\des données};
\node (prep) at (-30:3.0cm) {Préparation\\des données};
\node (model) at (-90:3.0cm) {Modélisation};
\node (eval) at (-150:3.0cm) {Évaluation};
\node (deploy) at (150:3.0cm) {Déploiement};
\foreach \a/\b in {business/data,data/prep,prep/model,model/eval,eval/deploy,deploy/business} {\draw[->,thick,ensiblue] (\a) to[bend left=12] (\b);}
\node[draw=none,fill=none,font=\normalsize\bfseries] at (0,0) {CRISP-DM\\adapté};
\end{tikzpicture}
\caption{Cycle CRISP-DM adapté au projet d'autoremplissage eCRF}
\label{fig:crispdm}
\end{figure}

\begin{table}[H]
\centering
\caption{Application de CRISP-DM au projet rapid-flow}
\begin{tabularx}{\textwidth}{p{0.28\textwidth}X}
\toprule
Phase CRISP-DM & Application au stage\\
\midrule
Compréhension métier & Identifier le besoin d'assistance à la saisie eCRF, les contraintes de traçabilité et le risque d'écriture non justifiée.\\
Compréhension des données & Étudier les bilans biologiques et comptes rendus d'imagerie, leurs unités, temporalités, libellés et variations de format PDF ou texte.\\
Préparation des données & Charger les documents, extraire le texte, reconstruire les lignes biologiques, structurer les sections et générer des chunks conservant le contexte.\\
Modélisation & Mettre en œuvre le routage documentaire, la recherche hybride, les extracteurs spécialisés, le mapping vers des clés canoniques et les règles de normalisation.\\
Évaluation & Comparer les prédictions aux gabarits de référence à l'aide de la précision, du rappel, du F1, de l'exactitude des valeurs et du faux autoremplissage.\\
Déploiement & Générer des exports JSON, CSV et XLS contrôlés, documenter l'exécution et intégrer les tests dans une chaîne d'intégration continue.\\
\bottomrule
\end{tabularx}
\end{table}

La méthodologie n'est pas strictement linéaire. Par exemple, les écarts constatés pendant l'évaluation de l'imagerie ont ramené le travail vers les phases de préparation et de modélisation : conservation du contexte complet, conversion cm/mm et priorité donnée à la conclusion RECIST. De même, les résultats biologiques ont motivé le renforcement de la reconstruction de lignes et du mapping d'alias. Cette boucle d'amélioration illustre l'intérêt de CRISP-DM pour conduire les itérations de manière contrôlée.

\section{Conclusion}
L'étude de l'existant confirme l'intérêt d'une chaîne hybride : règles et structures pour les données tabulaires, recherche ciblée pour réduire le contexte, extraction spécialisée pour les documents narratifs, et validation avant écriture. Le chapitre suivant traduit ces besoins en une conception logicielle.
